\section{Combining strategies and extracting dominions}
\label{sec:strategies}

We prove determinacy together with a quantitative guarantee for the
strategies of \Max.  This guarantee controls the loss incurred when a
strategy is interrupted and later restarted.  We then use it to show that
an opponent dominion contains a dominion outside an attractor of
the least priority.  That statement is the structural input to the
recursive algorithm.

A fixed strategy $\sigma$ of \Max\ from a vertex $v$ has the
\emph{uniform prefix property} if
\begin{equation}
  \label{eq:uniform-prefix}
  \text{for every }\epsilon>0\text{ there is }C_\epsilon\geq 0
  \text{ such that }\wt(\rho)\geq-\epsilon t-C_\epsilon
\end{equation}
for every $\sigma$-consistent path $\rho$ of $t\geq0$ edges starting at $v$.
The constant is independent of the path and its length.  In particular,
the strategy is fixed before $\epsilon$ is chosen.  Dividing the bound by
$t$ shows that every consistent infinite play has lower mean payoff at
least zero.  For any finite collection of strategies satisfying
\eqref{eq:uniform-prefix}, the constants $C_\epsilon$ can be chosen in
common by taking their maximum.

Each fixed component strategy is restarted with its initial memory.
Its prefix bounds therefore apply to the history since the restart,
with constants independent of the preceding play.

\subsection{Combination across an attractor}

When the least priority is odd, repeated visits to it already give \Min\
a win.  When it is even, \Max\ must also control the accumulated rewards.
The next construction inserts increasingly long portions of a
mean-payoff strategy between successive attempts to visit that priority.
The qualitative combination argument and its increasing-round
construction appear in
\cite[Lemma~5; full version, Lemma~7]{BouyerMarkeyOlschewskiUmmels2011}.
We give a proof that tracks a uniform bound for all interrupted prefixes.

\begin{lemma}[Strategy combination]
\label{lem:combine}
Let $H$ be a nonempty arena, and let $d\geq0$ be an integer such that
$p(v)\geq d$ for every $v\in H$.  Define
\begin{gather*}
 P=\begin{cases}\Max,&d\text{ even},\\\Min,&d\text{ odd},\end{cases}
 \qquad
 S=\{v\in H:p(v)=d\},\\
 A=\Attr_P^H(S),\qquad R=H\setminus A.
\end{gather*}
The set $S$ is allowed to be empty.
\begin{enumerate}
\item If $d$ is odd and \Min\ wins everywhere in the game on $R$, then
\Min\ wins everywhere in $H$.
\item If $d$ is even, \Max\ wins for mean payoff alone everywhere in $H$,
and from every vertex of $R$ \Max\ has a conjunction-winning strategy in
$R$ satisfying \eqref{eq:uniform-prefix}, then from every vertex of $H$
\Max\ has a conjunction-winning strategy satisfying
\eqref{eq:uniform-prefix}.
\end{enumerate}
\end{lemma}

\begin{proof}
Suppose first that $d$ is odd.  On entering $R$, \Min\ starts a winning
strategy for the game on $R$ and follows it until play leaves $R$.
The set $R$ is closed at \Min's vertices, so such an exit can only be
chosen by \Max.  In $A$, \Min\ follows an attracting strategy to $S$.
After reaching $S$, take one edge, with an arbitrary choice if the vertex
belongs to \Min, and resume the same rule.

If $S$ is visited infinitely often, its odd priority is the least
priority occurring infinitely often, and \Min\ wins.  Otherwise play
eventually stays in $R$: each entry into $A$ would force another visit to
$S$ in bounded time.  The final stay in $R$ follows one winning strategy,
so prefix independence again gives a win for \Min.  The one-edge step
after reaching $S$ ensures that successive applications of the rule
advance the play.

Now let $d$ be even.  By Lemma~\ref{lem:mp-prefix}, fix a mean-payoff
strategy on $H$ whose every finite prefix, from every start, has weight
at least $-B$.  For each $x\in R$, also fix one conjunction-winning
strategy $\sigma_x$ in $R$ satisfying \eqref{eq:uniform-prefix}.  These
choices will remain fixed throughout the construction.

Starting at any vertex of $H$, \Max\ plays rounds $j=1,2,\ldots$.
Round $j$ has the following parts.
\begin{enumerate}
\item Follow the mean-payoff strategy for exactly $j$ edges.
\item If the current vertex is some $x\in R$, start $\sigma_x$ and follow
it for as long as play stays in $R$.  If play never leaves $R$, this is
the remainder of the infinite play.
\item Once play is in $A$, follow an attracting strategy to $S$.
End the round upon reaching $S$, counting this as already accomplished
if its first vertex in this part lies in $S$.
\end{enumerate}
The next round starts at the vertex of $S$ just reached.  Since $R$ is
closed at \Max's vertices, any exit in the second part is chosen by
\Min.  If $S=\varnothing$, then $R=H$ and the second part of the first
round simply lasts forever.

If infinitely many rounds finish, they give infinitely many visits to
$S$ at distinct times: round $j$ contains at least $j$ edges.  The least
infinitely recurring priority is then the even number $d$.  If only
finitely many rounds finish, the unfinished round has an infinite
second part, since its first and third parts have bounded length.  Its
tail follows a conjunction-winning strategy in $R$.  Thus the constructed
strategy satisfies parity in either case.

It remains to prove the uniform prefix property; this will also establish
the mean-payoff condition.  Let $W$ bound the absolute edge weights in
$H$, and fix $\delta>0$.  Choose a common constant $C_\delta$ for the
finitely many strategies $\sigma_x$; use $C_\delta=0$ if $R$ is empty.
Any internal prefix of the second part, of length $t'$, has weight at
least $-\delta t'-C_\delta$.  This applies even if the ambient play
subsequently exits $R$: the internal prefix is a consistent finite path
in the total arena $R$.  The exiting edge is counted separately and
costs at most $W$.  The attraction then uses at most $|H|$ edges.
The four contributions are displayed in Figure~\ref{fig:round}.
Every prefix of a round, including a completed round, therefore has
weight at least
\begin{equation}
  \label{eq:round-prefix}
  -\delta t-K_\delta,
  \qquad
  K_\delta=B+C_\delta+(1+|H|)W,
\end{equation}
where $t$ is the length of that prefix.

If $q$ rounds have completed by time $T$, then
$T\geq 1+2+\cdots+q$.  Including a possible unfinished round, at most
$1+\sqrt{2T}$ rounds contribute to the $T$-edge prefix $\rho_T$.
Summing \eqref{eq:round-prefix} gives
\[
  \wt(\rho_T)\geq-\delta T-K_\delta(1+\sqrt{2T}).
\]
For a given $\epsilon>0$, set $\delta=\epsilon/2$ and
$K=K_{\epsilon/2}$.  The elementary inequality
\[
  K\sqrt{2T}\leq\frac{\epsilon T}{2}+\frac{K^2}{\epsilon}
\]
shows that \eqref{eq:uniform-prefix} holds with
$C_\epsilon=K+K^2/\epsilon$.  Only the analysis, not the strategy,
depends on $\epsilon$.
\end{proof}

\begin{figure}[htbp]
\centering
\begin{tikzpicture}[>=Latex,font=\small,
 block/.style={draw,rounded corners=2pt,minimum height=1.05cm,align=center},
 every path/.style={line width=.6pt}]
 \node[block,fill=blue!8,draw=blue!60!black,text width=2.6cm] (mp)
   {Mean-payoff policy\\exactly $j$ edges};
 \node[block,fill=green!7,draw=green!40!black,text width=3.1cm,
   right=.28cm of mp] (res)
   {Internal edges in $R$\\until an exit, if entered};
 \node[block,fill=red!6,draw=red!55!black,text width=1.55cm,
   right=.28cm of res] (exit)
   {Exit edge\\if any};
 \node[block,fill=orange!12,draw=orange!65!black,text width=2.65cm,
   right=.28cm of exit] (attr)
   {Attract to $S$\\at most $|H|$ edges};
 \draw[->] (mp)--(res);
 \draw[->] (res)--(exit);
 \draw[->] (exit)--(attr);
 \node[below=.22cm of mp,align=center] {$\ge-B$};
 \node[below=.22cm of res,align=center] {$\ge-\delta t'-C_\delta$};
 \node[below=.22cm of exit,align=center] {$\ge-W$};
 \node[below=.22cm of attr,align=center] {$\ge-|H|W$};
 \node[above=.15cm of res,align=center,font=\footnotesize]
   {An infinite stay gives a winning tail};
\end{tikzpicture}
\caption{A round in the even-priority construction of
Lemma~\ref{lem:combine}, with lower bounds for the weights of finite
prefixes of its parts. The $R$-part and its exit edge are omitted when the mean-payoff
part ends in $A$. Block widths are schematic. Every completed round
$j$ uses at least $j$ edges, which bounds the number of completed rounds
by time $T$ by $O(\sqrt T)$.}
\label{fig:round}
\end{figure}


\subsection{Determinacy with uniform prefix bounds}

We next establish that the strategies required on $R$ in the combination
lemma are always available at the winning vertices of \Max.  Proving
this assertion jointly with determinacy avoids any assumption that an
arbitrary winning strategy already has uniform prefix bounds.
The induction follows the attractor reduction in
\cite[Algorithm~3 and Appendix~C]{ChatterjeeHenzingerSvozil2017},
carrying the quantitative property through each removal.

\begin{theorem}
\label{thm:determinacy}
In every finite arena, each vertex is winning for exactly one of \Max\
and \Min.  From each vertex winning for \Max, there is a
conjunction-winning strategy satisfying \eqref{eq:uniform-prefix}.
\end{theorem}

\begin{proof}
We prove both assertions by induction on the number of vertices.  They
are vacuous in the empty arena.  Let $H$ be nonempty and assume the
assertions for all smaller arenas.

First observe how a nonempty set of already established winners reduces
the induction.  Suppose $Y\subseteq H$ is nonempty and every vertex of
$Y$ is winning for a player $J$, with strategies satisfying
\eqref{eq:uniform-prefix} if $J=\Max$.  Every vertex of
$Z=\Attr_J^H(Y)$ is also winning for $J$: attract to $Y$, then start a
winning strategy from the vertex reached.  Attraction lasts at most
$|H|$ edges.  If $J=\Max$, its bounded loss and the finite choice of
target strategies preserve \eqref{eq:uniform-prefix}.

Apply the induction hypothesis to the smaller arena $H\setminus Z$.
It is closed at $J$'s vertices.  A strategy there for the opponent of
$J$ can therefore keep play in this subarena and remains winning in
$H$.  A strategy there for $J$ is followed until any exit to $Z$, at
which point $J$ starts the winning strategy in $H$ associated with the
entry vertex.  There is at most one such switch: the new strategy is
continued on all subsequent histories, even if it later leaves $Z$.
Prefix independence proves that this strategy wins.

When $J=\Max$, the switch also preserves the quantitative property.
For each $\epsilon>0$, the internal prefix in $H\setminus Z$ and the
prefix after the switch have lower bounds of the form
$-\epsilon t_1-C_1$ and $-\epsilon t_2-C_2$.  The intervening edge costs
at most $W$, where $W$ bounds the absolute weights in $H$.  Taking the
constants in common over the finitely many possible entry vertices
gives a bound with slope $-\epsilon$ and constant $C_1+C_2+W$.
Prefixes before a switch satisfy the first bound.  If $J=\Min$, the
winning strategies of \Max\ remain entirely in the subarena and retain
their prefix bounds without switching.  Thus finding such a nonempty
$Y$ proves both assertions for $H$.

Let $d=\min_{v\in H}p(v)$, let $P$ be \Max\ for even $d$ and \Min\
for odd $d$, and let $Q$ be the other player.  Set
\[
 S=\{v\in H:p(v)=d\},\qquad
 A=\Attr_P^H(S),\qquad R=H\setminus A.
\]
If $d$ is even and \Max\ does not win for mean payoff alone everywhere
in $H$, Proposition~\ref{prop:mp-interface} supplies a nonempty set of
vertices from which \Min\ can force failure of the mean-payoff
condition.  These are conjunction-losing vertices, so they provide
$Y$ with $J=\Min$ in the preceding reduction.  We may therefore assume
that the mean-payoff hypothesis of Lemma~\ref{lem:combine} holds whenever
$d$ is even.

Since $S$ is nonempty, $R$ has fewer vertices than $H$.  Apply induction
to the game on $R$.  If its $Q$-winning set is nonempty, these strategies
remain winning in $H$: $R$ is closed at $P$'s vertices, so $Q$ can keep
play there.  Their prefix bounds also remain valid when $Q=\Max$.
They provide the required set $Y$ with $J=Q$.

In the remaining case, $P$ wins everywhere in $R$, with the uniform
prefix property when $P=\Max$.  Lemma~\ref{lem:combine} now proves the
assertions for all of $H$.  Finally, a vertex cannot be winning for both
players, since their objectives are complementary.
\end{proof}

\subsection{A dominion outside the least-priority attractor}

We can now turn the combination lemma around.  If an opponent dominion
had no opponent-winning part outside the relevant attractor, the lemma
would give the other player a winning strategy throughout that dominion.
For parity games this is the extraction principle of
\cite[Lemma~2.3]{LehtinenParysScheweWojtczak2022}; here the even case
also requires the ordinary mean-payoff hypothesis.

\begin{lemma}[Dominion extraction]
\label{lem:extract}
Let $H$ be a nonempty arena and $d\geq0$ an integer such that
$p(v)\geq d$ for all $v\in H$.  Let $P=\Max$ if $d$ is even and
$P=\Min$ if $d$ is odd, and let $Q$ be the other player.  Put
\[
 S=\{v\in H:p(v)=d\},\qquad
 R=H\setminus\Attr_P^H(S).
\]
If $d$ is even, assume that \Max\ wins for mean payoff alone everywhere
in $H$.  Then every nonempty $Q$-dominion $D$ in $H$ contains a nonempty
$Q$-dominion of the game on $R$.  The conclusion allows
$S=\varnothing$.
\end{lemma}

\begin{proof}
Inside the induced game on $D$, form
\[
  R_D=D\setminus\Attr_P^D(D\cap S).
\]
Because $D$ is a $Q$-dominion, it is closed at $P$'s vertices in $H$.
The attractor complement $R_D$ is closed at $P$'s vertices in $D$, hence
also in $H$.  It is disjoint from $S$, so the attractor-avoidance
observation in Section~\ref{sec:arenas} gives $R_D\subseteq R$.

If $d$ is even, \Max\ wins for mean payoff alone everywhere in $D$ as
well.  Indeed, the choices of \Max\ under a winning strategy in $H$ remain
available in $D$, because $D$ is closed at \Max's vertices; \Min\ is
restricted to choosing edges in $D$.

Suppose that the game on $R_D$ has no $Q$-winning vertex.
Theorem~\ref{thm:determinacy} makes every vertex there winning for $P$,
with the uniform prefix property if $P=\Max$.  Apply
Lemma~\ref{lem:combine} to the arena $D$, using the target $D\cap S$
and its residual arena $R_D$.  All its hypotheses have just been
verified, so $P$ wins everywhere in $D$, contrary to the fact that $D$
is a nonempty $Q$-dominion.  This argument includes $R_D=\varnothing$,
when its winning-strategy hypothesis is vacuous.

Consequently the $Q$-winning set in $R_D$ is nonempty.  It is a
$Q$-dominion there, and it lifts to a $Q$-dominion in $R$ by
Lemma~\ref{lem:arena-facts}, since $R_D$ is closed at $P$'s vertices.
It is contained in $D$, as required.  If $S=\varnothing$, then
$R=H$ and $R_D=D$; the same argument applies.
\end{proof}
